\section{Block Composition, Strong Powers, and Asymptotic Capacity}\label{sec:capacity-theory}
%==============================================================================


This section lifts the one-shot graph characterization to strong powers and asymptotic rates. The graph characterization already forces a block theory, because repeated composition preserves confusability structure rather than erasing it. The next results identify the correct strong-power object and the resulting asymptotic zero-error rate.

A genuine supermultiplicative block law holds:
\[
\alpha_{m+n}\;\ge\;\alpha_m\alpha_n,
\]
where $\alpha_n$ denotes the largest one-shot zero-error code size in the $n$-block composed system. This is the correct discrete precursor to Shannon-capacity analysis: the admissible code-size sequence is already supermultiplicative at the graph level.\leanmeta{\LH{MFT22}}

More sharply, the block-level graph itself now factors correctly. The mechanized development proves both a generic strong-power addition theorem for graphs and the induced model-specific statement that confusability between concatenated $(m+n)$-block states is equivalent to adjacency in the strong product of the $m$-block and $n$-block confusability graphs, packaged as an explicit graph isomorphism. Thus the repeated-block system is not only bounded by strong-product arguments at the codebook level; it is literally a strong-product graph-power object in the mechanized theory.\leanmeta{\LHrng{MFT}{29}{30}, \LHrng{MFT}{38}{39}}

Iterating the same construction yields the derived $n$-block lower bound. If $\alpha_1$ denotes the one-shot maximum independent-set code size, then the mechanized block-power theorem gives
\[
\alpha_1^n \;\le\; \alpha_n,
\]
where $\alpha_n$ is the largest one-shot zero-error code size in the $n$-fold block-composed system. This is the first asymptotic-capacity scaffold in the extended model: repeated block composition already forces the expected exponential growth law before any Shannon-capacity or $\vartheta$-style upper theory is introduced.\leanmeta{\LH{MFT21}}

For notation, write $\Theta(G)$ for the Shannon capacity of a confusability graph $G$ in logarithmic units, and write $\vartheta_\infty(G)$ for the corresponding asymptotic Lov\'asz-$\vartheta$ upper value. When the object is a view family rather than an abstract graph, we write $C_\infty(\text{views})$ and $\vartheta_\infty(\text{views})$ for the same quantities computed on its induced confusability graph.

These normalized block rates can be packaged into an explicit lower asymptotic capacity envelope
\[
C_* \;:=\; \sup_{n\ge 1} \frac{\log \alpha_n}{n},
\]
implemented as an $\mathrm{ENNReal}$ supremum of the positive block-rate sequence. In the mechanized development, $\alpha_n$ is now identified directly with the maximum finite independent-set size of the $n$-block confusability graph, the block graph is identified with the $n$-fold strong power of the one-shot confusability graph, and the same envelope is proved equal to the graph-side Shannon lower-capacity expression built from those strong powers.\leanmeta{\LHrng{MFT}{31}{32}, \LHrng{MFT}{35}{36}}

The asymptotic scaffold is also not limited to a lower envelope statement. On the graph side itself, strong-product lower bounds for independent-set cardinality induce monotonicity of normalized strong-power rates along repeated multiples. Specializing those generic graph theorems back to the model-generated confusability graph yields the repeated-block statement below.\leanmeta{\LHrng{MFT}{40}{42}}

Repeating a fixed $m$-block system $k$ times yields
\[
\alpha_m^k \;\le\; \alpha_{km},
\]
and hence
\[
\frac{\log \alpha_m}{m} \;\le\; \frac{\log \alpha_{km}}{km}.
\]
Thus rates along repeated block multiples are monotone from below toward the same capacity envelope. This is still a lower theory rather than a full Shannon-capacity theorem, but it is already a genuine asymptotic graph-rate statement inside the mechanized development.\leanmeta{\LH{MFT33}}

\begin{theorem}[Asymptotic Shannon-capacity theorem]\label{thm:asymptotic-capacity}
Let $\alpha_n$ denote the largest one-shot zero-error code size in the $n$-block composed system. Then the normalized block rates
\[
\frac{\log \alpha_n}{n}
\]
converge to a real asymptotic Shannon-capacity value, and that value is equal to the supremum of the finite block-rate sequence.
\end{theorem}
\leanmeta{\LHrng{MFT}{43}{49}, \LH{MFT56}}

\begin{proof}
The block law above gives supermultiplicativity of the code-size sequence, equivalently subadditivity of the negative log-independence sequence on strong powers. The trivial vertex-count bound supplies a lower bound on the normalized quotients, so Fekete's lemma applies and yields convergence of the normalized rates. The same mechanized argument then identifies the limit with the supremum of the finite block-rate sequence and matches the older $\mathrm{ENNReal}$ lower-capacity envelope with the resulting real asymptotic value.
\end{proof}

\paragraph{Relation to graph capacity theory.}
This places the extended model directly in the classical zero-error graph-capacity setting. The difference from the classical channel picture is generative rather than formal: here the confusability graphs arise from deterministic view families and partial observations on latent tuples. We do not claim a new general Shannon-capacity theorem for arbitrary graphs; the new content is that this partial-view model reaches the classical graph machinery in a non-clique regime and supports a structurally characterized equality subclass. The clique-fiber regime recovers the classical cluster-graph equality case, while the transitivity criterion identifies when that case is produced by the view-family model. The non-clique witness already shows that the model also generates genuinely nontrivial graph behavior.
